%% Section 8 -- Related Work
%% Ported from docs/DRAFT_PAPER.md Section 8.
%%
%% Four in-text tags in the Markdown draft were topic labels rather than
%% author-year (docs/DRAFT_PAPER.md drafting note 16). Converted here to real
%% citation keys, which resolves note 16 for the LaTeX source:
%%   [SoK 2025]          -> breitinger2025sok
%%   [SLR 2025]          -> onik2025slr
%%   [Resilience 2025]   -> wunder2025resilience
%%   [Preprocessing 2026]-> chaudhary2026preprocessing
%%
%% Section 8.2's final paragraph is CONTINGENT (docs/DRAFT_PAPER.md drafting
%% note 3). It is written so that narrowing removes a sentence rather than
%% restructuring the section. Keep it that way.

\section{Related Work}
\label{sec:related}

CARTA sits at the meeting point of four bodies of work, none of which addresses
the reconstruction of consumer chat archives after a destructive restore, and all
of which supply pieces of the machinery. The argument of this section is narrow
and worth stating up front. Merging chat archives is not an open problem; working
tools do it today. Staging forensic reconstruction as a pipeline is not new; it
was published in 2014. Per-record provenance is not new either; forensics
formalised it in the same year. What has not been specified is the contract such a
reconstruction owes its reader: what enters and leaves each stage, what becomes of
a record that cannot be safely interpreted, and what the procedure must report
about its own behaviour.

\subsection{Optimistic replication and replicated data types}
\label{sec:related:replication}

Divergent replicas are the founding concern of optimistic replication. Parker and
colleagues formalised the detection of mutual inconsistency between copies as
early as 1983~\citep{parker1983mutual}, and Coda supplies the closest structural
analogue to TCF's sequence: a client works disconnected, both sides advance, and
the system reintegrates on reconnection, naming and flagging the conflicts it
cannot resolve~\citep{kistler1992coda}. Coda attempts reintegration. Condition 3
describes a platform that does not.

Convergent and commutative replicated data types give conditions under which
divergent replicas provably converge. It is worth being precise about what those
conditions are, because the obvious reading of them is wrong for our purposes.
Shapiro and colleagues define state-based merge as a least upper bound over a join
semilattice, taking two states and no ancestor at all~\citep{shapiro2011crdt}.
Ancestry enters elsewhere: the mergeable replicated data types line models state
as versions that ``trace their origin to a common ancestor version'' and merges via
three-way merge against a lowest common ancestor~\citep{kaki2019mergeable}, and the
certified treatment makes the existence of that ancestor a store precondition,
assuming execution begins with a single branch so that any two branches have an
LCA~\citep{soundarapandian2022certified}. Reconciliation literatures beyond CRDTs
share the assumption. Log-based reconciliation takes a common initial state plus
per-replica action logs, and the formal treatment of filesystem synchronisation
requires the update sequence to be applicable to the original filesystem
state~\citep{csirmaz2022sync}.

What excludes TCF from the CRDT guarantee is therefore Condition 2, not ancestry.
Convergence in the state-based model holds on the assumption that the system
transmits payloads between replicas infinitely often over eventually reliable
channels, and only so long as all replicas receive all
updates~\citep{shapiro2011crdt}. TCF's branches exchange nothing during the
divergence window. Nothing prevents a message set from being treated as a join
semilattice under union after the fact; the difficulty is historical rather than
algebraic, in that no synchronising representation existed while the two branches
were diverging, so no convergence guarantee was ever purchased. Kleppmann and
colleagues put the design premise directly, describing CRDTs as multi-user from the
ground up, and note that side-by-side document versions and branches remain an
area with little work behind it~\citep{kleppmann2019localfirst}.

Retrofitting reconciliation onto systems not built for it has been addressed in one
direction only. Conflict-free replicated relations augment conventional relational
schemas with causal metadata so that existing applications gain multi-synchronous
access with little modification~\citep{yu2020crr}. That retrofit is prospective:
the metadata is in place before divergence begins. CARTA addresses the
retrospective case, in which the artefacts already exist and no such metadata was
ever recorded. \citet{csirmaz2022sync} characterise why this regime resists the OT
and CRDT framings, contrasting their constant low-latency communication and small
deltas with synchronisation that happens in a single exchange after structural
differences have accumulated over an extended period. TCF is that regime without
the original state their model still assumes.

\subsection{Timeline reconstruction in digital forensics}
\label{sec:related:forensics}

Forensic timeline reconstruction is the closest neighbour by problem shape, and the
field's own systematisation describes a four-stage process of extraction,
normalisation, arrangement and analysis, with multi-source combination handled
informally~\citep{breitinger2025sok}. Provenance, quarantine, conflict detection,
determinism and invariants do not appear there as formally defined terms.

Chabot and colleagues supply the most developed formal apparatus in this lineage,
and two of their contributions bound our claim. The knowledge-representation paper
defines a relation of support linking a footprint to the entity it was used to
deduce, carried through the analysis and materialised in the output fact
base~\citep{chabot2014knowledge}. Per-record provenance is thus formalised in
forensics, and we do not claim otherwise. The relation is an
identifier-to-identifier derivation edge; CARTA's provenance additionally binds a
byte range and a content hash to each record and states an invariant requiring
every synthesised record to carry at least one provenance entry
(Section~\ref{sec:carta}). The architecture paper published the same year defines a
four-layer pipeline running from extraction through semantic interpretation and
reasoning to an investigator-facing interface, covering the path ``from the
extraction of data to the analysis of the timeline
produced''~\citep{chabot2014timeline}. Staging a reconstruction process is
therefore established prior art, and CARTA claims no priority over it.

The distinctions that remain are specific. That architecture describes its
inter-stage relationship in prose, saying that normalisation translates
extraction-layer data into the format the upper layer uses, without naming or
typing that format; no stage carries a signature or a pre- and post-condition, and
the completeness requirement is attributed to earlier work rather than discharged.
Records that are irrelevant or unusable are filtered out to reduce downstream
volume, which is the opposite disposition from quarantine. The correlation operator
quantifies relatedness between two events in order to surface investigative links,
and does not collapse two observations of one event into a canonical record; we
flag the resemblance to CARTA's similarity-based deduplication explicitly because
the two operate on similar inputs toward opposite ends. The outputs are a queryable
triple store and a visualisation interface rather than a generated report.

\subsection{Timeline tooling and chat-archive tooling}
\label{sec:related:tooling}

Plaso, with its predecessor log2timeline, is the reference open-source
super-timeline engine, originally presented as an automation framework for
producing a combined timeline from heterogeneous
sources~\citep{gudjonsson2010log2timeline}. Its treatment of failure is more
developed than a summary dismissal would suggest: it defines five typed warning
containers covering analysis, extraction, preprocessing, recovery and timelining,
serialises them into the storage file, exposes them through \texttt{pinfo}, and
flags event data in-band as corrupted or recovered. Records that fail are retained
and inspectable. What they are not is part of the timeline output contract, since
the export path iterates sorted events alone, so the artefact an analyst reads
carries no representation of what failed. Plaso's documentation states no
determinism, correctness or completeness property, and describes its own event
model as currently free form with formal specification tracked as open work.

Its multi-source behaviour is concatenation followed by exact-duplicate
suppression, and the suppression key incorporates the source data stream's path
specification, so the same event observed in two independent sources is by
construction never collapsed into one record. Conflicting observations across
sources are not detected. Suppressed events are reported as an aggregate count,
without recording which events were dropped or which surviving record absorbed
them.

Chat-specific forensic work is abundant and single-acquisition. Anglano's analysis
of WhatsApp on Android is the canonical reference, and its reconstruction claim is
explicitly the correlation of artefacts within one device to recover contacts and
message chronology~\citep{anglano2014whatsapp}; a recent systematic review of
secure instant-messaging forensics shows a taxonomy organised around acquisition,
decryption and artefact interpretation, with no category for reconciling multiple
archives~\citep{onik2025slr}.

The strongest prior art is not in the literature at all. Working tools merge chat
archives today. \texttt{bepaald/\allowbreak signalbackup-tools}, at 1,392 stars and under
active development, imports one Signal backup into
another~\citep{repo_signalbackuptools}; its documentation declines to deduplicate,
advising users whose backups partly overlap to crop the source by date before
merging. \texttt{natario1/\allowbreak whatsapp-database-merger} merges decrypted WhatsApp
message databases with attention to referential integrity, deduplicating only
incidentally through database uniqueness constraints, and drops media it cannot
copy~\citep{repo_wadbmerger}. Smaller projects come closer to CARTA's stated
target: \texttt{tarikab7/\allowbreak whatsapp-chat-merger} merges plain-text WhatsApp exports
into a unified timeline with claimed automatic deduplication and per-source colour
tagging~\citep{repo_wachatmerger}, and \texttt{frozenspider/\allowbreak chat-history-manager}
normalises Telegram, WhatsApp, Signal and several other formats into one schema,
whose author notes candidly that duplicate removal may also discard legitimate
repeated messages~\citep{repo_chathistorymanager}. \texttt{EkriirkE/\allowbreak WAMerge} exists
solely because its author encountered the TCF sequence and wrote a tool for
it~\citep{repo_wamerge}. Commercial products advertise the same capability.

These tools bound the contribution in a way the academic record does not. Merging
is available, the problem is recognised, and neither can be claimed. What they
establish instead is the gap. The tools differ in how they treat input they cannot
parse, from silent discard through a printed per-record diagnostic to aborting the
load outright, and none of them retains rejected input as an inspectable output.
None couples merging with per-record provenance:
\texttt{SecurityRonin/\allowbreak chat4n6} tags every recovered record with its evidence source
and a confidence score but works on one acquisition and performs no
merge~\citep{repo_chat4n6}, while the tools that merge carry no provenance. No
merging tool found reports conflicts or states a determinism property, and the one
deduplication predicate we examined in source ---
\texttt{tarikab7/\allowbreak whatsapp-chat-merger}'s --- keys on timestamps floored into
two-second buckets, which makes detection depend on where a pair falls relative to
a bucket boundary rather than on how far apart the pair is. Adjacent formats have
better-developed deduplication semantics, notably
\texttt{ragibson/\allowbreak SMS-MMS-deduplication}, which enumerates attachment presence,
country-code variance, timestamp precision and malformed encodings while taking
care to avoid false positives, and which a reviewer may reasonably read as the
closest existing analogue to CARTA's deduplication
criterion~\citep{repo_smsmmsdedup}.

Work of similar shape is appearing in adjacent domains. A deterministic
preprocessing framework for heterogeneous network forensic datasets addresses
incompatible schemas and timestamp formats and offers a mechanism for verifying
consistency across runs~\citep{chaudhary2026preprocessing}. The problem statement
is close enough that the difference is worth stating rather than assuming: that
work canonicalises datasets for correlation, while CARTA specifies reconciliation
of overlapping archives of one conversation, with quarantine and reporting as part
of the output contract.

\subsection{Personal information management and multi-device practice}
\label{sec:related:pim}

Personal archiving research describes the conditions in which TCF becomes possible.
\citet{marshall2008archiving} observes that individuals' digital belongings are
distributed across many stores, some overlapping and some independent, and
identifies device replacement as a point where local files are lost for good. That
divergence is derivational, between edited variants of one artefact, and the loss
is genuine destruction, so it is a different configuration from the one defined
here.

Two survey studies support the motivation without evidencing the phenomenon. A
three-country survey of 1,423 participants found that 46\% had experienced at least
one data-loss incident and that, among those who recovered using backups, over half
had backups that were outdated or incomplete~\citep{wunder2025resilience};
incomplete restores are the common case rather than the exception. An earlier panel
study reports that fewer than half of young adults back up regularly, though it
models backup frequency as a protective behaviour rather than studying device
replacement~\citep{redmiles2019phone}. Neither describes two independently
persistent archives of one account, and we have found no HCI or CSCW study that
does.
